\subsection{N-011: why the named obstructions occur at the Laurent edges}
\label{sec:nonstandard-edges}

The inherited negative-coordinate and trace obstructions are not being
claimed again.  The following elementary transfer identity explains
their locations and distinguishes the two different edge patterns.
Let
\[
 f(D)=\sum_{j=0}^m a_jD^j,\quad a_m\ne0,\qquad
 g(v)=\delta^\varepsilon f(D),\quad
 \delta=v+v^{-1},\quad D=\delta^2,\quad\varepsilon\in\{0,1\}.
\]
Put $L=2m+\varepsilon$.  Since
$\delta^n=v^n(1+v^{-2})^n$, the coefficient of $v^{L-2r}$ is
\[
 [v^{L-2r}]g(v)=
 \sum_{k=0}^{\min(r,m)}
 a_{m-k}\binom{L-2k}{r-k},
 \tag{N-33}
\]
where an out-of-range binomial coefficient is zero.  In particular
the outermost and next coefficients are
\[
             a_m,\qquad La_m+a_{m-1}.
 \tag{N-34}
\]
The coefficients at the other end are equal by bar invariance.
Thus the highest $D$-term alone controls the extreme Laurent degree;
lower $D$-degrees cannot repair its sign.  At the next Laurent degree,
the second-highest $D$-term competes with the binomial contribution
from the highest one.

For the exact examples requested in item 7, this gives:
\[
\begin{array}{l|r|r|r|r|r}
\text{example}&L&a_m&a_{m-1}&[v^L]&[v^{L-2}]\\\hline
\text{GCT\_NSRH F57}&18&-1&20&-1&2\\
\text{INT-NSRH-4: }(2013,3201)&15&-5&102&-5&27\\
\text{INT-NSRH-4: }(2013,3021)&16&-3&84&-3&36\\
\text{INT-NSRH-4: }(2103,3021)&15&-3&70&-3&25\\
\text{F67, three-plus-three}&12&2&-36&2&-12\\
\text{F67, one-plus-three}&8&2&-20&2&-4
\end{array}
 \tag{N-35}
\]
The three INT-NSRH-4 pairs use that intake's zero-based permutation
notation; transposing each pair supplies its other three coordinates.
The exact full $D$-polynomials and Laurent lists are in
\nolinkurl{edge_coefficients_verified.json}.  For example the first is
\[
\delta\bigl(-5D^7+102D^6-515D^5+1856D^4-3888D^3
                     +5568D^2-4608D+2048\bigr).
\]
All three have a negative highest $D$-coefficient, followed by a
sufficiently large positive next coefficient.  Their negative Laurent
coefficients are consequently at the outermost positions; the exact
remaining coefficient lists certify positivity inward from there.
For F57 the inherited polynomial starts $-D^9+20D^8-\cdots$, giving
the same mechanism.  This is compatible with the direct mixed-descent
calculation in N-002: virtual negative terms can dominate the highest
degree while positive lower-degree terms contribute only farther in.

F67 is a different pattern: its outermost coefficients are positive,
and its negative coefficients occur one step inward.  For the
three-plus-three example, the inherited direct four-dimensional
tensor calculation gives
\[
 \tau_6(D)=D^6+(-D^3+18D^2-48D+32)^2
          =2D^6-36D^5+\cdots.
\]
The negative cross term in the square outweighs the binomial spill
from $2D^6$: $-36+12\cdot2=-12$.  For the one-plus-three supplement,
\[
 \tau_{1,3}(D)=2D^4-20D^3+84D^2-128D+64,
\]
and the analogous coefficient is $-20+8\cdot2=-4$.
The source tensor rule includes the complementary, virtual generators
$mI-N_i$.  Its whole-product trace has those negative cross terms
despite the input matrices being effective.  A positive value of the
sum-of-squares trace at a real parameter does not imply coefficientwise
positivity.

These explanations concern the named inherited examples.  They prove
no theorem that every negative coordinate or word trace must have its
negative coefficients at the edges.  In particular, F67 does not
support a claim that the outermost coefficient itself must be negative.

\paragraph{Provenance and exact checks.}
The inputs are GCT\_NSRH F57; INT-NSRH-4 in
\nolinkurl{INTAKE_GCT_NSRH_20261121_part3.md}, sections 1--2;
and GCT\_NSRH F67 with its one-plus-three supplement.  The script
\nolinkurl{edge_coefficients.py} recomputes the three ambient coordinates
to obtain their new $D$-expressions, verifies that they and their
transposes are precisely the six bad coordinates of the given word,
and checks (N-33) at every exponent of all six displayed cases.
For F67 it also derives the scalar polynomials directly from the four
by four output matrices.  The representation-theoretic occurrence
claims and original obstructions retain their inherited finding IDs.
